\begin{abstract}
Agentic Reinforcement Learning（RL）使大语言模型（LLM）能够执行自主决策与长程规划。与传统 LLM 后训练不同，agentic RL 的工作负载高度异构，同时包含计算密集的 prefill、带宽受限的 decoding，以及具备状态、且严重依赖 CPU 的环境模拟。我们认为，要高效训练 agentic RL，必须采用解耦式基础设施，以便针对不同子任务使用最合适的硬件；但若直接解耦，又会因为各阶段间依赖复杂而带来显著的同步开销和资源闲置。

本文提出 \SystemName{}，这是一个面向解耦式基础设施、用于最大化多任务 agentic RL 吞吐的分布式系统。\SystemName{}建立在三个核心原则之上：1）\emph{硬件亲和的工作负载映射}，将计算受限与带宽受限任务路由到最匹配的 GPU；2）\emph{细粒度异步}，以轨迹为粒度组织执行，减少流水线气泡；3）\emph{状态感知计算}，把无状态组件（如奖励模型）卸载到无服务器平台以获得弹性伸缩。实验表明，\SystemName{}能够显著提升训练吞吐，并相较单体式和同步式基线将端到端训练时间缩短 1.35--2.05 倍。我们还在阿里巴巴超过 3000 张 GPU 的集群上，用其训练了一个百亿级以上参数规模的 MoE 模型用于 Qoder 产品，进一步验证了系统的可扩展性与鲁棒性。代码地址：\url{https://github.com/alibaba/ROLL}。
\end{abstract}
